\section*{अध्याय \ref{ch:b3:computational-chemistry} --- अभिकलनात्मक रसायन: हार्ट्री--फ़ॉक और DFT}

\begin{solution}{exo:b3:computational-chemistry:1}
$\dd E/\dd\alpha = \frac32 - \sqrt{2/\pi\alpha} = 0$ से $\alpha = 8/9\pi = 0.283\,
a_0^{-2}$ और $E = \frac32\alpha - 2\sqrt{2\alpha/\pi} = -4/3\pi = \qty{-0.4244}{\hartree}$,
जो यथार्थ $-0.5$ से 15\,\% ऊपर है: एक अकेले गाउसी फलन में $1s$ फलन का न नुकीला शीर्ष है
न पूँछ।
\end{solution}

\begin{solution}{exo:b3:computational-chemistry:2}
STO-3G: प्रति C 5 फलन ($1s$, $2s$, तीन $2p$) और प्रति H 1: $6 \times 5 + 6 =
36$। 6-31G(d): प्रति C, 1 (कोर) + $2 \times 4$ (विभक्त संयोजकता) + 6 ($d$) $= 15$;
प्रति H, 2: $6 \times 15 + 6 \times 2 = 102$।
\end{solution}

\begin{solution}{exo:b3:computational-chemistry:3}
इलेक्ट्रॉनिक \omterm{def:b3:quantum-model-systems:schrodinger}{हैमिल्टनी संकारक} में नाभिकीय द्रव्यमान नहीं होते: \omterm{def:b3:computational-chemistry:born-oppenheimer}{बॉर्न--ओपेनहाइमर
सन्निकटन} के भीतर PES, अतः $r_e$ और $k = U''(r_e)$ भी,
समान हैं। $\tilde\omega = \sqrt{k/\mu}/2\pi c$ \omterm{def:b3:quantum-model-systems:oscillator}{समानीत द्रव्यमान} पर निर्भर है, जो
दुगुना हो जाता है: अनुपात $\sqrt2 = 1.414$ (मापित 1.413)।
\end{solution}

\begin{solution}{exo:b3:computational-chemistry:4}
आठ परमाणु $3N - 6 = 18$ कंपन देते हैं: 17 वास्तविक और एक काल्पनिक। हेसियन का एक
ऋणात्मक \omterm{def:b3:quantum-model-systems:operator}{अभिलक्षणिक मान}: प्रथम-कोटि पल्याण बिंदु, एक संक्रमण
संरचना (यहाँ, उदाहरण के लिए, किसी घूर्णन या किसी विलोपन की)।
\end{solution}

\begin{solution}{exo:b3:computational-chemistry:5}
$\zeta = 27/16$, $E = -(27/16)^2 = \qty{-2.8477}{\hartree} = \qty{-77.49}{eV}$।
त्रुटि $-2.8477 + 2.9034 = \qty{0.0557}{\hartree} = \qty{1.52}{eV}$, अर्थात्
कुल ऊर्जा का 1.9\,\% --- परंतु अनेक अभिक्रिया ऊर्जाओं से बड़ी।
\end{solution}

\begin{solution}{exo:b3:computational-chemistry:6}
$(-1 - E)(-0.5 - E) - (-0.2 - 0.3E)^2 = 0$, अर्थात् $0.91E^2 + 1.38E + 0.46 = 0$:
$E = \qty{-1.0217}{\hartree}$ और \qty{-0.4947}{\hartree}। हाँ: दूसरे फलन को मिलाने से
ऊर्जा $H_{11}$ से नीचे आ जाती है, जैसा \omterm{thm:b3:computational-chemistry:variational}{विचरण सिद्धांत}
अनुमति देता है।
\end{solution}

\begin{solution}{exo:b3:computational-chemistry:7}
$(102/36)^4 = 64$: लगभग 640 मिनट, कोई 11 घंटे।
\end{solution}

\begin{solution}{exo:b3:computational-chemistry:8}
$0.5782 \times 27.211 = \qty{15.73}{eV}$, मापे गए मान से 0.30 eV ऊपर। स्थिर
कक्षक: वास्तविक आयन शिथिल होता है और नीचे होता है, अतः कूपमान्स का मान बहुत ऊँचा है।
लुप्त सहसंबंध: उदासीन अणु (दो इलेक्ट्रॉन) में आयन (एक इलेक्ट्रॉन) से अधिक \omterm{def:b3:computational-chemistry:correlation}{सहसंबंध
ऊर्जा} है, जो वास्तविक मान को बड़ा बनाती है।
यहाँ पहली त्रुटि भारी पड़ती है।
\end{solution}

\begin{solution}{exo:b3:computational-chemistry:9}
$-0.5960 - 2(-0.4666) = \qty{0.3372}{\hartree} = \qty{9.18}{eV}$, दो परमाणुओं से ऊपर।
RHF फलन में 50\,\% \ce{H+ H-} बना रहता है, और एक प्रोटॉन तथा एक
हाइड्राइड को अलग करने में H की आयनन ऊर्जा और H की इलेक्ट्रॉन बंधुता का
अंतर लगता है, जिसे ऊर्जा औसत में सम्मिलित कर लेती है।
\end{solution}

\begin{solution}{exo:b3:computational-chemistry:10}
$E(c_1^2 + 2c_1c_2S + c_2^2) = c_1^2H_{11} + 2c_1c_2H_{12} + c_2^2H_{22}$।
$c_1$ के सापेक्ष अवकलन करके, $\partial E/\partial c_1 = 0$ के साथ:
$E(2c_1 + 2c_2S) = 2c_1H_{11} + 2c_2H_{12}$, अर्थात् $(H_{11} - E)c_1 + (H_{12} -
ES)c_2 = 0$; इसी प्रकार $(H_{12} - ES)c_1 + (H_{22} - E)c_2 = 0$।
\end{solution}

\begin{solution}{exo:b3:computational-chemistry:11}
$k \approx [E(1.30) - 2E(1.35) + E(1.40)]/(0.05)^2 = 0.001417/0.0025 =
\qty{0.567}{\hartree\per\bohr\squared} = \qty{882}{N/m}$। $\mu = m_H/2 =
\qty{8.37e-28}{kg}$ के साथ $\omega = \sqrt{k/\mu} = \qty{1.03e15}{s^{-1}}$ और
$\tilde\omega = \omega/2\pi c = \qty{5452}{cm^{-1}}$, जो मापे गए
$\tilde\omega_e$ से 24\,\% ऊपर है: न्यूनतम-आधार RHF वक्र बहुत खड़ा है (सहसंबंध नहीं, आधार
बहुत छोटा), इसीलिए परिकलित आवृत्तियों को घटाकर मापित किया जाता है।
\end{solution}

\begin{solution}{exo:b3:computational-chemistry:12}
हाइड्रोजन आबंध वाले कुछ \unit{kJ/mol} के अंतरों के लिए सहसंबंध और
परिक्षेपण चाहिए: परिक्षेपण संशोधन के साथ एक संकर फलनक, या अंतिम ऊर्जाओं के लिए
बेहतर एक सहसंबद्ध तरंग-फलन विधि, \omterm{def:b3:computational-chemistry:split-valence}{विसरित फलनों} सहित ध्रुवित
त्रि-विभक्त आधार के साथ। दोनों संरूपणों को
एक ही स्तर पर इष्टतम कीजिए, आवृत्तियों से निम्निष्ठ की पुष्टि कीजिए (सभी वास्तविक), \omterm{def:b3:quantum-model-systems:oscillator}{शून्य-बिंदु
ऊर्जाएँ} जोड़िए, एक बड़े परिकलन से आधार को परखिए, और ज्ञात संरूपण ऊर्जा वाले
किसी संबंधित निकाय से तुलना कीजिए।
\end{solution}

\begin{solution}{pb:b3:computational-chemistry:1}
\textbf{1.} न्यूनतम आधार के हर \omterm{def:b3:computational-chemistry:basis}{स्लेटर-प्रकार कक्षक} के स्थान पर उस पर समंजित तीन
गाउसी फलनों का एक नियत संकुचन रखा जाता है।
\textbf{2.} हीलियम नाभिक अपने इलेक्ट्रॉनों को अधिक आकर्षित करता है, अतः उसका $1s$ फलन
अधिक सघन है (बड़े घातांक)। $6.3624/3.4253 = 1.857 = (1.69/1.24)^2$।
\textbf{3.} दो आधार फलन, दो इलेक्ट्रॉन, एक अधिकृत कक्षक (और एक
रिक्त)।
\textbf{4.} $V_{\mathrm{nn}} = Z_{\mathrm{He}}Z_{\mathrm H}/R = 2/1.4632 =
\qty{1.3669}{\hartree}$।
\textbf{5.} दोनों फलन प्रबलता से अतिव्यापन करते हैं (54\,\%): परमाणु इतने
निकट हैं कि आबंध बना सकें।
\textbf{6.} $h_{11}$, He फलन में एक इलेक्ट्रॉन की ऊर्जा है, दोनों
नाभिकों के साथ परंतु किसी अन्य इलेक्ट्रॉन के बिना; He नाभिक (आवेश 2) उसे कहीं अधिक
दृढ़ता से थामता है।
\textbf{7.} $(h_{11} - \varepsilon)(h_{22} - \varepsilon) - (h_{12} - \varepsilon S)^2 = 0$:
$0.71185\varepsilon^2 + 2.79292\varepsilon + 2.44969 = 0$।
\textbf{8.} $\varepsilon = \qty{-2.600}{\hartree}$ और \qty{-1.324}{\hartree}।
\textbf{9.} $\mathbf h$ दोनों इलेक्ट्रॉनों के बीच के प्रतिकर्षण की उपेक्षा करता है; कोर
कक्षक बहुत संकुचित और बहुत नीचा है।
\textbf{10.} $F_{\mu\nu} = h_{\mu\nu} + \sum_{\lambda\sigma}P_{\lambda\sigma}[(\mu\nu|\sigma\lambda) -
\frac12(\mu\lambda|\sigma\nu)]$, जहाँ $P_{\lambda\sigma} = 2C_{\lambda1}C_{\sigma1}$।
\textbf{11.} हर इलेक्ट्रॉन का दूसरे के आवेश-मेघ द्वारा कूलॉम प्रतिकर्षण (और,
सामान्यतः, विनिमय)।
\textbf{12.} $\mathbf F$, $\mathbf P$ पर निर्भर है, जो
$\mathbf F$ से प्राप्त कक्षकों पर निर्भर है: समीकरण अरैखिक हैं और उत्तरोत्तर
सन्निकटनों द्वारा स्व-संगति तक हल किए जाते हैं।
\textbf{13.} $c_1^2 + c_2^2 + 2c_1c_2S = 0.7684 + 0.0410 + 0.1906 = 1.0000$।
\textbf{14.} He: $1.5369 + 0.1906 = 1.727$; H: $0.0820 + 0.1906 = 0.273$
इलेक्ट्रॉन।
\textbf{15.} आवेश: He $+0.27$, H $+0.73$। धनावेश अधिकतर
हाइड्रोजन पर है: \ce{HeH+} एक हीलियम परमाणु से बँधा प्रोटॉन है, \ce{He + H+}, जैसा
अपेक्षित है, क्योंकि He की आयनन ऊर्जा (24.6 eV) H की (13.6 eV) से कहीं अधिक है।
\textbf{16.} $-\varepsilon_1 = \qty{1.633}{\hartree} = \qty{44.4}{eV}$।
\textbf{17.} $E_{\mathrm{el}} = -2.8418 - 1.3669 = \qty{-4.2087}{\hartree}$।
\textbf{18.} तीन चक्र \qty{-2.8418}{\hartree} तक पहुँचते हैं। हर चक्र एक
सारणिक देता है, जिसकी ऊर्जा अभिसरित हार्ट्री--फ़ॉक
ऊर्जा की उपरि सीमा है (\omterm{thm:b3:computational-chemistry:variational}{विचरण सिद्धांत}), और पुनरावृत्तियाँ उसे सुधारती हैं।
\textbf{19.} $\zeta = 2.0925$ वाला आधार निम्नतर ऊर्जा देता है, अतः इस अणु के लिए वही
बेहतर है (\omterm{thm:b3:computational-chemistry:variational}{विचरण सिद्धांत}): हीलियम फलन
\ce{HeH+} में मुक्त परमाणु की तुलना में अधिक सघन है, जिसके $\zeta = 1.69$ की
मानक आधार नकल करता है।
\textbf{20.} रिक्त प्रतिआबंधी कक्षक $\sigma^*$; उसकी ऊर्जा \ce{HeH+} की इलेक्ट्रॉन
बंधुता के ऋणात्मक का सन्निकटन होती (इतने छोटे आधार में ख़राब ढंग से)।
\textbf{21.} शून्य (कोई इलेक्ट्रॉन नहीं); $\qty{-2.8078}{\hartree}$।
\textbf{22.} $-2.8078 - (-2.8418) = \qty{0.0340}{\hartree} = \qty{89}{kJ/mol}$।
\textbf{23.} मापे गए मान का लगभग आधा: न्यूनतम आधार हीलियम घनत्व को
प्रोटॉन की ओर ध्रुवित नहीं कर सकता (कोई $p$ फलन नहीं हैं), अतः आबंध
बहुत दुर्बल है। (शून्य-बिंदु और \qty{298}{K} संशोधन इसकी तुलना में छोटे हैं।)
\textbf{24.} $-(24.5874 + 54.4178)/27.211 = \qty{-2.9034}{\hartree}$; STO-3G परमाणु
\qty{0.0956}{\hartree} (\qty{2.6}{eV}) बहुत ऊँचा है।
\textbf{25.} ज्यामितियाँ इस पर निर्भर करती हैं कि ऊर्जा $R$ के साथ कैसे बदलती है; अधिकांश
त्रुटि, जो कोर में केंद्रित है, हर ज्यामिति पर लगभग समान है और
निरस्त हो जाती है।
\textbf{26.} \textbf{RHF/STO-3G (मानक आधार) पर, \qty{1.4632}{\bohr} पर,
$E(\ce{HeH+}) = \qty{-2.8418}{\hartree}$} (और चिरप्रचलित
$\zeta_{\mathrm{He}} = 2.0925$ के साथ \qty{-2.8607}{\hartree})।
\end{solution}
